Activation steering---adding a learned direction to a language model's hidden states at inference time---is among the simplest methods for controlling model behavior without fine-tuning. Practitioners must choose whether to steer continuously (risking degenerate output) or for a limited window, yet the actual dynamics of how steering influence persists after the vector is removed have never been directly measured. We present the first systematic characterization of steering vector decay in transformer hidden states. Using contrastive activation addition for sycophancy on \qwen (28 layers, 1.5B parameters), we apply the steering vector for $N \in \{1, 3, 5, 10\}$ initial tokens, then measure the signed projection of the hidden-state perturbation onto the steering direction at each subsequent step under teacher-forced generation. We find that steering influence decays to less than 5\% of its peak within a single token after the vector is removed, following an exponential with time constants $\tau \approx 1.2$--$5.8$ tokens ($R^2$ up to $0.97$). The KV cache is the primary vehicle for any residual signal: preserving it retains a small positive trace, while resetting it produces a statistically significant negative rebound (overcorrection; Cohen's $d$ up to $3.41$, $p < 0.0001$). These results demonstrate that continuous steering is necessary because transformer dynamics aggressively wash out transient perturbations, and they provide mechanistic support for KV-cache-based steering as a more persistent alternative.
